\chapter{Julia Implementation Reference}
\label{app:julia}

% ============================================================================
% Appendix B: Julia Code Reference
% Written: 2025-12-31
% ============================================================================

This appendix provides Julia equivalents for all estimators presented in the main text. The Julia implementations follow the SciML Problem-Estimator-Solution pattern and achieve numerical parity with the Python implementations to 10 decimal places.

\section{Installation}

\begin{verbatim}
cd julia
julia --project -e "using Pkg; Pkg.instantiate()"
\end{verbatim}

\section{Architecture Overview}

The Julia implementation follows the SciML pattern:

\begin{enumerate}
    \item \textbf{Problem}: Immutable struct containing data and parameters
    \item \textbf{Estimator}: Algorithm choice (e.g., \texttt{SimpleATE}, \texttt{DoublyRobust})
    \item \textbf{Solution}: Results struct with estimates, SEs, CIs
\end{enumerate}

\begin{minted}{julia}
# General pattern
problem = SomeProblem(outcomes, treatment, covariates, parameters)
solution = solve(problem, SomeEstimator())

# Access results
solution.estimate   # Point estimate
solution.se         # Standard error
solution.ci_lower   # Lower CI bound
solution.ci_upper   # Upper CI bound
\end{minted}

% ============================================================================
\section{RCT Estimators}
\label{sec:julia_rct}
% ============================================================================

\subsection{Simple ATE (Difference-in-Means)}

\begin{minted}{julia}
"""
    SimpleATE <: AbstractRCTEstimator

Simple difference-in-means with Neyman heteroskedasticity-robust variance.

    tau_hat = Y_bar_1 - Y_bar_0
    SE(tau_hat) = sqrt(s1^2/n1 + s0^2/n0)
"""
struct SimpleATE <: AbstractRCTEstimator end

function solve(problem::RCTProblem, estimator::SimpleATE)::RCTSolution
    (; outcomes, treatment, parameters) = problem

    y1 = outcomes[treatment]
    y0 = outcomes[.!treatment]
    n1, n0 = length(y1), length(y0)

    # Point estimate
    ate = mean(y1) - mean(y0)

    # Neyman variance
    se = sqrt(var(y1)/n1 + var(y0)/n0)

    # Confidence interval (Satterthwaite df)
    df = satterthwaite_df(var(y1), n1, var(y0), n0)
    ci_lower, ci_upper = confidence_interval_t(ate, se, parameters.alpha, df)

    return RCTSolution(estimate=ate, se=se, ci_lower=ci_lower, ...)
end
\end{minted}

\textbf{Usage:}
\begin{minted}{julia}
using CausalEstimators

outcomes = [10.0, 12.0, 4.0, 5.0]
treatment = [true, true, false, false]
problem = RCTProblem(outcomes, treatment, nothing, nothing, (alpha=0.05,))
solution = solve(problem, SimpleATE())
\end{minted}

% ============================================================================
\section{Propensity Score Methods}
\label{sec:julia_psm}
% ============================================================================

\subsection{Propensity Score Estimation}

\begin{minted}{julia}
"""
    estimate_propensity_scores(treatment, covariates) -> PropensityResult

Estimate propensity scores via logistic regression.
"""
function estimate_propensity_scores(treatment, covariates)
    X = hcat(ones(size(covariates, 1)), covariates)
    coef = logistic_regression(X, treatment)
    propensity = logistic.(X * coef)
    auc = compute_propensity_auc(propensity, treatment)
    return PropensityResult(propensity, auc, coef)
end
\end{minted}

\subsection{Nearest Neighbor Matching}

\begin{minted}{julia}
"""
    nearest_neighbor_matching(treated_ps, control_ps; with_replacement=true)

Match treated units to controls based on propensity score distance.
Returns indices of matched controls for each treated unit.
"""
function nearest_neighbor_matching(
    treated_ps::Vector{T},
    control_ps::Vector{T};
    with_replacement::Bool = true,
    caliper::Union{Nothing, T} = nothing,
) where {T<:Real}
    n_treated = length(treated_ps)
    matches = Vector{Int}(undef, n_treated)

    available = trues(length(control_ps))

    for i in 1:n_treated
        distances = abs.(control_ps .- treated_ps[i])

        if !with_replacement
            distances[.!available] .= Inf
        end

        best_idx = argmin(distances)

        # Check caliper
        if !isnothing(caliper) && distances[best_idx] > caliper
            matches[i] = 0  # No valid match
        else
            matches[i] = best_idx
            available[best_idx] = false
        end
    end

    return matches
end
\end{minted}

% ============================================================================
\section{Weighting Estimators}
\label{sec:julia_weighting}
% ============================================================================

\subsection{Inverse Probability Weighting (IPW)}

\begin{minted}{julia}
"""
    ObservationalIPW <: AbstractObservationalEstimator

IPW estimator using Hajek (normalized) weights.

    tau_IPW = [sum T_i*Y_i/e(X_i)] / [sum T_i/e(X_i)] -
              [sum (1-T_i)*Y_i/(1-e(X_i))] / [sum (1-T_i)/(1-e(X_i))]
"""
struct ObservationalIPW <: AbstractObservationalEstimator end

function solve(problem::ObservationalProblem, ::ObservationalIPW)
    (; outcomes, treatment, covariates) = problem

    # Estimate propensity scores
    prop_result = estimate_propensity_scores(treatment, covariates)
    e = prop_result.propensity

    # Compute IPW weights
    weights = ifelse.(treatment, 1 ./ e, 1 ./ (1 .- e))

    # Hajek estimator (normalized)
    mu1 = sum(weights[treatment] .* outcomes[treatment]) / sum(weights[treatment])
    mu0 = sum(weights[.!treatment] .* outcomes[.!treatment]) / sum(weights[.!treatment])
    ate = mu1 - mu0

    # Influence function variance
    influence = @. ifelse(treatment, (outcomes - mu1)/e, -(outcomes - mu0)/(1-e))
    se = sqrt(var(influence) / length(outcomes))

    return IPWSolution(estimate=ate, se=se, ...)
end
\end{minted}

\subsection{Doubly Robust (AIPW)}

\begin{minted}{julia}
"""
    DoublyRobust <: AbstractObservationalEstimator

AIPW estimator combining IPW with outcome regression.

Double robustness: consistent if EITHER propensity OR outcome model correct.

    tau_DR = (1/n) sum_i [T_i/e(X_i)*(Y_i - mu1(X_i)) + mu1(X_i)
                        - (1-T_i)/(1-e(X_i))*(Y_i - mu0(X_i)) - mu0(X_i)]
"""
struct DoublyRobust <: AbstractObservationalEstimator end

function solve(problem::ObservationalProblem, ::DoublyRobust)
    (; outcomes, treatment, covariates) = problem
    T = convert(Vector{Float64}, treatment)

    # Estimate propensity and outcome models
    e = estimate_propensity_scores(treatment, covariates).propensity
    mu0, mu1 = fit_outcome_models(outcomes, treatment, covariates)

    # AIPW formula
    treated_contrib = @. T/e * (outcomes - mu1) + mu1
    control_contrib = @. (1-T)/(1-e) * (outcomes - mu0) + mu0

    ate = mean(treated_contrib .- control_contrib)

    # Influence function variance
    influence = treated_contrib .- control_contrib .- ate
    se = sqrt(mean(influence.^2) / length(outcomes))

    return DRSolution(estimate=ate, se=se, ...)
end
\end{minted}

% ============================================================================
\section{CATE Estimators}
\label{sec:julia_cate}
% ============================================================================

\subsection{S-Learner}

\begin{minted}{julia}
"""
    SLearner <: AbstractCATEEstimator

Single model approach: fit mu(X, T) -> Y, then CATE(x) = mu_hat(x, 1) - mu_hat(x, 0).
"""
struct SLearner <: AbstractCATEEstimator end

function solve(problem::CATEProblem, ::SLearner)
    (; outcomes, treatment, covariates) = problem
    n = length(outcomes)

    # Augment covariates with treatment
    X_aug = hcat(ones(n), covariates, Float64.(treatment))
    coef = fit_linear_model(X_aug, outcomes)

    # CATE = coefficient on treatment (for linear model)
    # For flexibility, compute explicitly:
    X_treat = hcat(ones(n), covariates, ones(n))
    X_ctrl = hcat(ones(n), covariates, zeros(n))

    cate = X_treat * coef .- X_ctrl * coef  # = coef[end] for all i
    ate = mean(cate)

    return CATESolution(cate=cate, ate=ate, ...)
end
\end{minted}

\subsection{T-Learner}

\begin{minted}{julia}
"""
    TLearner <: AbstractCATEEstimator

Two model approach: fit mu1(X) on treated, mu0(X) on control.
CATE(x) = mu1_hat(x) - mu0_hat(x)
"""
struct TLearner <: AbstractCATEEstimator end

function solve(problem::CATEProblem, ::TLearner)
    (; outcomes, treatment, covariates) = problem

    # Fit separate models
    X1 = hcat(ones(sum(treatment)), covariates[treatment, :])
    X0 = hcat(ones(sum(.!treatment)), covariates[.!treatment, :])

    coef1 = fit_linear_model(X1, outcomes[treatment])
    coef0 = fit_linear_model(X0, outcomes[.!treatment])

    # Predict for all units
    X_all = hcat(ones(length(outcomes)), covariates)
    cate = X_all * coef1 .- X_all * coef0

    return CATESolution(cate=cate, ate=mean(cate), ...)
end
\end{minted}

% ============================================================================
\section{Causal Forests}
\label{sec:julia_forests}
% ============================================================================

Causal forests are implemented via \texttt{econml} in Python. For Julia usage, call Python via \texttt{PyCall.jl}:

\begin{minted}{julia}
using PyCall

econml = pyimport("econml.dml")

function causal_forest(outcomes, treatment, covariates; n_estimators=100, honest=true)
    model = econml.CausalForestDML(
        n_estimators=n_estimators,
        honest=honest,
        cv=5,
    )
    model.fit(Y=outcomes, T=treatment, X=covariates)

    cate = model.effect(covariates)
    ate_inf = model.ate_inference(X=covariates)

    return (
        cate = vec(cate),
        ate = mean(cate),
        ate_se = ate_inf.stderr_mean,
    )
end
\end{minted}

\begin{insight}
Critical: Always use \texttt{honest=true} for valid confidence intervals. Without honest splitting, CIs have severe undercoverage (e.g., 72\% instead of 95\%). See CONCERN-28.
\end{insight}

% ============================================================================
\section{Sensitivity Analysis}
\label{sec:julia_sensitivity}
% ============================================================================

\subsection{E-Value}

\begin{minted}{julia}
"""
    compute_e_value(rr::Real) -> Float64

E-value: minimum confounding strength to explain away effect.

For RR >= 1: E = RR + sqrt(RR * (RR - 1))
For RR < 1: Apply formula to 1/RR
"""
function compute_e_value(rr::Real)::Float64
    rr > 0 || throw(ArgumentError("RR must be positive"))

    rr_use = rr < 1.0 ? 1.0/rr : rr

    if isapprox(rr_use, 1.0)
        return 1.0  # No effect
    end

    return rr_use + sqrt(rr_use * (rr_use - 1.0))
end

# Effect type conversions
smd_to_rr(d) = exp(0.91 * d)  # VanderWeele approximation
ate_to_rr(ate, p0) = (p0 + ate) / p0
\end{minted}

\textbf{Usage:}
\begin{minted}{julia}
problem = EValueProblem(2.5; ci_lower=1.8, ci_upper=3.5, effect_type=:rr)
solution = solve(problem, EValue())

println("E-value: $(solution.e_value)")
println("E-value for CI: $(solution.e_value_ci)")
\end{minted}

\subsection{Rosenbaum Bounds}

\begin{minted}{julia}
"""
    solve(problem::RosenbaumProblem, ::RosenbaumBounds)

Sensitivity analysis for matched pairs using Wilcoxon signed-rank test.

Finds critical Gamma where p_upper > alpha.
"""
function solve(problem::RosenbaumProblem, ::RosenbaumBounds)
    differences = problem.treated_outcomes .- problem.control_outcomes

    # Wilcoxon signed-rank statistic
    T_plus, ranks, signs = compute_signed_rank_statistic(differences)

    gamma_grid = range(problem.gamma_range..., length=problem.n_gamma)
    p_upper = zeros(length(gamma_grid))

    for (i, gamma) in enumerate(gamma_grid)
        E_upper, _, Var_max, _ = compute_bounds_at_gamma(ranks, signs, gamma)
        p_upper[i] = normal_approximation_p(T_plus, E_upper, Var_max)
    end

    # Critical Gamma: first where p_upper > alpha
    gamma_critical = findfirst(>(problem.alpha), p_upper)
    gamma_critical = isnothing(gamma_critical) ? nothing : gamma_grid[gamma_critical]

    return RosenbaumSolution(gamma_critical=gamma_critical, ...)
end
\end{minted}

\textbf{Usage:}
\begin{minted}{julia}
treated = [10.5, 12.3, 8.7, 15.2, 11.1]
control = [8.2, 10.1, 7.5, 12.8, 9.3]
problem = RosenbaumProblem(treated, control; gamma_range=(1.0, 3.0))
solution = solve(problem, RosenbaumBounds())

println("Critical Gamma: $(solution.gamma_critical)")
\end{minted}

% ============================================================================
\section{Cross-Language Validation}
\label{sec:julia_validation}
% ============================================================================

All Julia implementations are validated against Python to 10 decimal places:

\begin{minted}{julia}
# Test parity with Python
@test isapprox(julia_ate, python_ate, rtol=1e-10)
@test isapprox(julia_se, python_se, rtol=1e-10)
\end{minted}

Current validation status:
\begin{itemize}
    \item \textbf{RCT estimators}: 100\% parity
    \item \textbf{PSM}: 100\% parity
    \item \textbf{Weighting}: 100\% parity (IPW, DR, TMLE)
    \item \textbf{CATE}: 100\% parity (S/T/X/R-learners, DML)
    \item \textbf{Sensitivity}: 100\% parity (E-value, Rosenbaum)
\end{itemize}

